% concatenated from three_distributive_not_omega_distributive.tex
% Standalone source. Compile with pdflatex twice, or with latexmk -pdf.
\documentclass[11pt]{amsart}
\usepackage[T1]{fontenc}
\usepackage{lmodern}
\usepackage[a4paper,margin=28mm]{geometry}
\usepackage{amsmath,amssymb,amsthm,mathtools}
\usepackage{microtype}
\usepackage{enumitem}
\usepackage{needspace}
\usepackage[hidelinks]{hyperref}
\hypersetup{pdftitle={A continuous 3-distributive frame that is not omega-distributive},
  pdfsubject={A negative answer to Erne's distributivity and web-space questions},
  pdfkeywords={kappa-distributivity, wide coframe, continuous frame, wide web space}}

\newtheorem{theorem}{Theorem}[section]
\newtheorem{proposition}[theorem]{Proposition}
\newtheorem{lemma}[theorem]{Lemma}
\newtheorem{corollary}[theorem]{Corollary}
\theoremstyle{definition}
\newtheorem{definition}[theorem]{Definition}
\theoremstyle{remark}
\newtheorem{remark}[theorem]{Remark}
\numberwithin{equation}{section}
\newcommand{\NN}{\mathbb N}
\newcommand{\OO}{\mathcal O}
\newcommand{\PP}{\mathcal P}
\newcommand{\RR}{\overline{\mathbb R}_{+}}
\newcommand{\dn}{\mathord\downarrow}
\newcommand{\up}{\mathord\uparrow}
\newcommand{\DD}{\mathord\Downarrow}
\newcommand{\intr}{\operatorname{int}}
\newcommand{\meet}{\mathbin\wedge}
\newcommand{\tops}{\mathbf 1_X}
\newcommand{\bots}{\mathbf 0_X}
\setlist[enumerate]{label=\textup{(\roman*)},leftmargin=2.2em}

\title[Finite distributivity versus $\omega$-distributivity]
{A continuous $3$-distributive frame\break that is not $\omega$-distributive}
\date{}
\subjclass[2020]{06D10, 06B35, 54D45, 54H12}
\keywords{$\kappa$-distributivity, wide coframe, continuous frame,
wide web space, topological semilattice}

\author{Wei Luan}
\address{Key Laboratory of Computing and Stochastic Mathematics (Ministry of Education), School of Mathematics and Statistics, Hunan Normal University, Changsha, Hunan 410081, China}
\email{luanwei@hunnu.edu.cn}

\author{Qingguo Li}
\address{School of Mathematics, Hunan University, Changsha, Hunan, 410082, China}
\email{liqingguoli@aliyun.com}
	% Corresponding author}

\begin{document}
\begin{abstract}
We give a negative answer to the question, posed by Ern\'e, whether
every $3$-distributive lattice is $\omega$-distributive. More precisely,
we exhibit a continuous frame that is $\kappa$-distributive for every
integer $\kappa\geq 2$, but is not a wide coframe. The frame is the
open-set lattice of a compact, locally compact, countably based
$T_0$ topological meet-semilattice, obtained from Lawson's construction
in the logarithmic form described by Goubault-Larrecq. The failure of
$\omega$-distributivity is witnessed by an explicit matrix with
countably many nonempty finite rows: all row joins are the same nonzero
element, whereas every choice of one entry from each row has meet zero.
The same space answers negatively Ern\'e's accompanying question whether
every $4$-web space is a wide web space. All properties of the
construction needed for these conclusions are proved directly.
\end{abstract}
\maketitle

\section{Introduction}

Ern\'e \cite[Section~6]{Erne2009} introduced $\kappa$-distributivity
using covers of cardinality strictly smaller than $\kappa$.
For complete lattices, $\omega$-distributivity is equivalent to being
a \emph{wide coframe}, also called $\mathcal F$-distributivity.
After observing that $3$-distributivity implies $\kappa$-distributivity
for each finite $\kappa\geq 3$, he asked whether it also implies
$\omega$-distributivity \cite[p.~2065]{Erne2009}.
On the same page, he asked whether every $4$-web space is an
$\omega$-web space, that is, a wide web space.

We answer both questions negatively. The following statement uses
exactly the cardinal convention of \cite{Erne2009}.

\begin{theorem}\label{thm:main}
There exists a compact, locally compact, countably based $T_0$ space
$X$ whose specialization order has binary meets, with continuous
meet operation, such that its open-set lattice $L=\OO(X)$ satisfies:
\begin{enumerate}
\item $L$ is a continuous frame;
\item $L$ is $\kappa$-distributive for every integer $\kappa\geq 2$;
\item $L$ is not $\omega$-distributive, and hence is neither a wide
coframe nor completely distributive.
\end{enumerate}
Moreover, $X$ is a $\kappa$-web space for every integer $\kappa\geq 2$,
but is not a wide web space.
\end{theorem}

The underlying semilattice construction is not new. Lawson
\cite[Section~2, Example~1]{Lawson1970} constructed a compact Hausdorff
topological semilattice without a neighborhood base of subsemilattices
at its greatest element. We use the explicit logarithmic reformulation
and the non-Hausdorff topology described by Goubault-Larrecq
\cite{GL2022}. Our purpose is to derive the distributivity conclusions
above from this construction. We give the estimates, topological
arguments, and distributivity proofs in full; in particular, none of
the conclusions depends on an unproved transfer between Hausdorff,
Scott, and subspace topologies. The obstruction to
$\omega$-distributivity is presented as a countable matrix.

We work in ZFC. We write $\NN=\{1,2,\ldots\}$ and
$\NN_0=\{0,1,2,\ldots\}$. Empty joins and meets in a complete lattice
are its least and greatest elements, respectively. General background
on continuous lattices and Scott topology can be found in
\cite{GierzEtAl2003}.

\section{Distributivity and finite meet operations}\label{sec:prelim}

For a subset $A$ of a poset $P$, put
\[
\dn A=\{x\in P:(\exists a\in A)\ x\leq a\},\qquad
\up A=\{x\in P:(\exists a\in A)\ a\leq x\}.
\]
We write $\dn a$ and $\up a$ for singleton-generated sets.

\begin{definition}[{\cite[Section~6]{Erne2009}}]\label{def:kappa}
Let $\kappa\geq 2$ be a cardinal, or let $\kappa=\infty$.
Write
\[
\PP_\kappa L=\{A\subseteq L:|A|<\kappa\},\qquad
\PP_\infty L=\PP L.
\]
For $b$ in a complete lattice $L$, define
\begin{equation}\label{eq:approx}
\DD_\kappa b
 =\bigcap\left\{\dn A:A\in\PP_\kappa L,
                          \ b\leq\bigvee A\right\}.
\end{equation}
The lattice $L$ is \emph{$\kappa$-distributive} if
$b=\bigvee\DD_\kappa b$ for every $b\in L$.
\end{definition}

In particular, $\PP_\omega L$ consists of all finite subsets of $L$,
including the empty set. Consequently $\DD_\kappa 0=\varnothing$,
and the required equality at $0$ holds automatically. For $b\ne 0$,
membership $a\in\DD_3 b$ says precisely that
\[
b\leq c\vee d\quad\Longrightarrow\quad a\leq c\text{ or }a\leq d
\qquad(c,d\in L).
\]
Thus a convention using covers by $n$ elements corresponds to
Ern\'e's index $n+1$, with the empty-cover convention handled separately.

We record the equivalence with distributive identities, including a
proof to fix all conventions. Given a family $\mathcal A$ of subsets
of $L$, its \emph{crosscut system} is
\[
\mathcal A^{\#}
 =\left\{Z\subseteq\bigcup\mathcal A:
                 Z\cap A\ne\varnothing\text{ for every }A\in\mathcal A\right\}.
\]

\begin{lemma}\label{lem:distribution}
For a complete lattice $L$ and $\kappa$ as in
Definition~\ref{def:kappa}, the following are equivalent:
\begin{enumerate}
\item $L$ is $\kappa$-distributive;
\item for every $\mathcal A\subseteq\PP_\kappa L$,
\begin{equation}\label{eq:crosscut}
\bigwedge_{A\in\mathcal A}\bigvee A
 =\bigvee_{Z\in\mathcal A^{\#}}\bigwedge Z;
\end{equation}
\item for every set $I$, sets $J_i$ with $|J_i|<\kappa$, and
elements $a_{i,j}\in L$, one has
\begin{equation}\label{eq:matrix}
\bigwedge_{i\in I}\bigvee_{j\in J_i}a_{i,j}
 =\bigvee_{f\in\prod_{i\in I}J_i}\bigwedge_{i\in I}a_{i,f(i)}.
\end{equation}
For $\kappa=\infty$, there is no restriction on the sets $J_i$.
\end{enumerate}
\end{lemma}

\begin{proof}
The right-hand side of \eqref{eq:crosscut} is always at most its
left-hand side. Assume (i), and write the left-hand side as $b$.
For $a\in\DD_\kappa b$, the set
$Z_a=(\bigcup\mathcal A)\cap\up a$ meets every $A\in\mathcal A$,
by \eqref{eq:approx}. Thus $Z_a\in\mathcal A^\#$ and
$a\leq\bigwedge Z_a$. Taking the join over $\DD_\kappa b$ proves (ii).
This also covers an empty $\mathcal A$, with $Z_a=\varnothing$.

Conversely, assume (ii) and fix $b\in L$. Set
\[
\mathcal A_b=\{A\in\PP_\kappa L:b\leq\bigvee A\}.
\]
Since $\{b\}\in\mathcal A_b$, the left-hand side of
\eqref{eq:crosscut} for this family is $b$.
For each $Z\in\mathcal A_b^\#$ and $A\in\mathcal A_b$, an element
$z\in Z\cap A$ satisfies $\bigwedge Z\leq z$. Hence
$\bigwedge Z\in\DD_\kappa b$. Formula~\eqref{eq:crosscut} gives
$b\leq\bigvee\DD_\kappa b$; the reverse inequality follows by
using $A=\{b\}$ in \eqref{eq:approx}.

Finally, the equivalence of (ii) and (iii) follows from the Axiom of
Choice. For clarity, if all rows are nonempty, the set of selected
values $\{a_{i,f(i)}:i\in I\}$ is a crosscut of the family of row
sets. Conversely, for each crosscut $Z$, choose in each row a value
belonging to $Z$; the meet of those selected values is at least
$\bigwedge Z$. If a row is empty, both sides of
\eqref{eq:matrix} are $0$. If $I=\varnothing$, both sides are $1$.
Taking the rows to be arbitrary members of $\mathcal A$ proves
the converse implication as well.
\end{proof}

The terminology in \cite[Sections~1 and~6]{Erne2009} therefore makes
$\omega$-distributivity exactly the wide coframe property, and
$\infty$-distributivity exactly complete distributivity in ZFC.
The coframe property is equivalent to \emph{fixed $3$-distributivity},
which requires \eqref{eq:crosscut} for
$\mathcal A\subseteq\PP_3L$ with $\bigcap\mathcal A\ne\varnothing$.
We will establish the full condition of Definition~\ref{def:kappa}.

For a space $Y$, its specialization preorder is defined by
\[
x\leq y\quad\Longleftrightarrow\quad
\text{every open set containing $x$ contains $y$}.
\]
All open sets are upper sets for this preorder. In a $T_0$ space it is
an order. The complete lattice $\OO(Y)$ has joins and meets
\begin{equation}\label{eq:opens}
\bigvee\mathcal U=\bigcup\mathcal U,
\qquad
\bigwedge\mathcal U=\intr_Y\left(\bigcap\mathcal U\right).
\end{equation}

\begin{lemma}\label{lem:finite}
Let $Y$ be a $T_0$ space whose specialization order has binary meets.
If the meet map $Y\times Y\longrightarrow Y$ is continuous, then
$\OO(Y)$ is $\kappa$-distributive for every integer $\kappa\geq 2$.
\end{lemma}

\begin{proof}
Fix $\kappa\geq 2$ and set $n=\kappa-1$. The map
\[
\mu_n:Y^n\longrightarrow Y,
\qquad \mu_n(y_1,\ldots,y_n)=y_1\wedge\cdots\wedge y_n,
\]
is continuous, by induction, with $\mu_1$ the identity.
Let $U$ be a nonempty open set and $x\in U$. Continuity at
$(x,\ldots,x)$ gives open neighborhoods $V_1,\ldots,V_n$ of $x$
such that $\mu_n(\prod_{i=1}^n V_i)\subseteq U$.
For $V=\bigcap_{i=1}^n V_i$, one has
\begin{equation}\label{eq:finite-local}
x\in V\subseteq U,\qquad \mu_n(V^n)\subseteq U.
\end{equation}

We claim that $V\in\DD_\kappa U$ in $\OO(Y)$.
Any family of fewer than $\kappa$ open sets covering $U$ is nonempty
and can be listed, with repetitions, as $U_1,\ldots,U_n$.
If $V\nsubseteq U_i$ for every $i$, choose $x_i\in V\setminus U_i$.
Then $z=\bigwedge_{i=1}^n x_i\in U$, so $z\in U_j$ for some $j$.
Since $z\leq x_j$ and $U_j$ is upper, this implies $x_j\in U_j$,
a contradiction. The claim follows from \eqref{eq:approx}.
Every $x\in U$ belongs to such a $V$, proving
$U=\bigvee\DD_\kappa U$. For $U=\varnothing$, the approximation
set is empty and the equality is immediate.
\end{proof}

\section{The semilattice construction}\label{sec:construction}

We now specify the version of the Lawson construction used below;
the logarithmic formula is from \cite{GL2022}.
For real $t\geq 1$, define
\[
h(t)=\log_2\bigl(\log_2(1+t)\bigr).
\]
The function $h$ is increasing, $h(1)=0$, and
\begin{equation}\label{eq:slow}
h(2k)-h(k)
 =\log_2\!\left(\frac{\log_2(1+2k)}{\log_2(1+k)}\right)
 \longrightarrow 0\qquad(k\longrightarrow\infty).
\end{equation}
For $i\in\NN_0$, put
\begin{equation}\label{eq:coordinates}
M_i=2^{2^i}-1,\qquad N_i=\{1,\ldots,M_i\},\qquad S_i=\PP(N_i).
\end{equation}
Order $S_i$ by inclusion, and define $\sigma_i:S_i\to[0,\infty]$ by
\begin{equation}\label{eq:sigma}
\sigma_i(A)=
\begin{cases}
\infty,&A=N_i,\\
i-h(|N_i\setminus A|),&A\ne N_i.
\end{cases}
\end{equation}
Each $\sigma_i$ is monotone. Since $h(M_i)=i$, for proper $A$ we have
$0\leq\sigma_i(A)\leq i$, and in particular
\begin{equation}\label{eq:sigma-special}
\sigma_i(\varnothing)=0,
\qquad \sigma_i(N_i\setminus\{k\})=i\quad(k\in N_i).
\end{equation}

\begin{lemma}\label{lem:slow-intersection}
For real numbers $\tau>\varepsilon>0$, there is $i_0\in\NN_0$ such
that for all $i\geq i_0$ and $A,B\in S_i$,
\[
\sigma_i(A)>\tau,\quad\sigma_i(B)>\tau
\quad\Longrightarrow\quad
\sigma_i(A\cap B)>\tau-\varepsilon.
\]
\end{lemma}

\begin{proof}
By \eqref{eq:slow}, choose an integer $K\geq1$ such that
$h(2k)<h(k)+\varepsilon$ for every integer $k\geq K$.
Choose $i_0>\tau-\varepsilon+h(2K)$. Fix $i\geq i_0$.
If $A=N_i$ or $B=N_i$, the conclusion follows immediately.
Otherwise, set
\[
a=|N_i\setminus A|,\quad b=|N_i\setminus B|,\quad
k=\max\{a,b\},\quad c=|N_i\setminus(A\cap B)|.
\]
Then $1\leq c\leq a+b\leq 2k$. If $k\geq K$, then
\[
i-h(c)>i-h(k)-\varepsilon
 =\min\{\sigma_i(A),\sigma_i(B)\}-\varepsilon
 >\tau-\varepsilon.
\]
If $k<K$, monotonicity of $h$ gives
\[
i-h(c)\geq i-h(2K)\geq i_0-h(2K)>\tau-\varepsilon.
\]
In both cases, $i-h(c)=\sigma_i(A\cap B)$.
\end{proof}

Let $\RR=[0,\infty]$ with its usual order and Scott topology.
The nonempty proper open sets of $\RR$ are the intervals
$(r,\infty]$ with $0\leq r<\infty$. Give each finite lattice $S_i$
its upper-set topology, which equals its Scott topology. Let
\begin{equation}\label{eq:ambient}
K=\RR\times\prod_{i\in\NN_0}S_i
\end{equation}
with the product topology, and give
\begin{equation}\label{eq:X}
X=\{(s,(A_i)_{i\in\NN_0})\in K:
                         s\leq\sigma_i(A_i)\text{ for all }i\in\NN_0\}
\end{equation}
the subspace topology. Throughout the proof this is the specified
topology on $X$; it is not replaced by a Scott topology on $X$.
A base consists of
\begin{equation}\label{eq:base}
B(J;m,E)=X\cap
 \left(J\times\prod_{i=0}^{m}\up E_i\times\prod_{i>m}S_i\right),
\end{equation}
where $m\in\NN_0$, $E_i\subseteq N_i$, and either $J=\RR$ or
$J=(r,\infty]$ for a real $r\geq0$.

\begin{proposition}\label{prop:order}
The specialization order of $X$ is the componentwise order:
\[
(s,(A_i))\leq(t,(B_i))
\quad\Longleftrightarrow\quad s\leq t\text{ and }A_i\subseteq B_i
\text{ for every }i.
\]
In particular, $X$ is $T_0$. Its least and greatest elements are
\[
\bots=(0,(\varnothing)_i),\qquad
\tops=(\infty,(N_i)_i).
\]
Every nonempty open subset of $X$ contains $\tops$.
\end{proposition}

\begin{proof}
The basic opens in \eqref{eq:base} are upper sets for the
componentwise order, proving one implication. Conversely, if $s>t$,
then $t<\infty$ and the open cylinder with first coordinate in
$(t,\infty]$ contains $(s,(A_i))$ but not $(t,(B_i))$.
If $A_j\nsubseteq B_j$, choose $e\in A_j\setminus B_j$; the open
cylinder requiring $e$ in the $j$th coordinate separates the two
points in the same direction. This proves the order assertion and
the $T_0$ property. The formulas for the bounds follow from
\eqref{eq:sigma-special} and $\sigma_i(N_i)=\infty$.
The last assertion follows because opens are upper sets.
\end{proof}

\begin{proposition}\label{prop:meet}
For $x=(s,(A_i))$ and $y=(t,(B_i))$ in $X$, their greatest lower
bound is
\begin{equation}\label{eq:meet}
x\wedge y=(u,(A_i\cap B_i)),\qquad
u=\min\left\{s,t,\inf_{i\in\NN_0}\sigma_i(A_i\cap B_i)\right\}.
\end{equation}
\end{proposition}

\begin{proof}
The point on the right belongs to $X$ and is below $x$ and $y$.
If $(v,(D_i))\in X$ is another common lower bound, then
$v\leq s,t$ and $D_i\subseteq A_i\cap B_i$ for every $i$.
Monotonicity of $\sigma_i$ yields
$v\leq\sigma_i(D_i)\leq\sigma_i(A_i\cap B_i)$ for every $i$.
Thus $v\leq u$, proving the assertion.
\end{proof}

\section{Compactness and continuity}\label{sec:topology}

Give $\RR$ its compact order topology, denoted by $\RR^{\sharp}$,
and each $S_i$ its discrete topology, denoted by $S_i^{\sharp}$.
The map $s\mapsto s/(1+s)$, with $\infty\mapsto1$, identifies
$\RR^{\sharp}$ with $[0,1]$. The product
\[
K^{\sharp}=\RR^{\sharp}\times\prod_{i\in\NN_0}S_i^{\sharp}
\]
is compact Hausdorff by the Tychonoff theorem, and its topology
refines the topology of $K$.

\Needspace{8\baselineskip}
\begin{lemma}\label{lem:closed}
The set $X$ is closed in $K^{\sharp}$.
\end{lemma}

\begin{proof}
For every $i$, the subset
\[
G_i=\bigcup_{A\in S_i}[0,\sigma_i(A)]\times\{A\}
      \subseteq\RR^{\sharp}\times S_i^{\sharp}
\]
is closed, being a finite union of closed sets.
The constraint $s\leq\sigma_i(A_i)$ is the inverse image of $G_i$
under the corresponding coordinate projection. The intersection
of these closed constraints is $X$.
\end{proof}

\begin{proposition}\label{prop:compact}
The space $X$ is compact and countably based. For each $x\in U\in\OO(X)$,
there is a compact saturated set $C$ such that
\begin{equation}\label{eq:compact-nbhd}
x\in\intr_X C\subseteq C\subseteq U.
\end{equation}
In particular, $X$ is locally compact.
\end{proposition}

\begin{proof}
Lemma~\ref{lem:closed} and the coarseness of the topology of $X$
give compactness. In \eqref{eq:base} one may restrict $r$ to
nonnegative rational numbers. There are countably many resulting
basic opens, since the coordinate spaces are finite.

Choose $B(J;m,E)$ with $x\in B(J;m,E)\subseteq U$.
If $J=\RR$, put $C=B(J;m,E)$. This set is closed in $K^{\sharp}$:
the conditions $A_i\supseteq E_i$ constrain finitely many finite
discrete coordinates. It is therefore compact in $X$, and it is open.
If $J=(r,\infty]$ and the first coordinate of $x$ is $s$, choose a
finite real number $t$ with $r<t<s$. Such a choice also exists when
$s=\infty$. Set
\[
C=X\cap\left([t,\infty]\times\prod_{i=0}^{m}\up E_i
                               \times\prod_{i>m}S_i\right).
\]
This set is closed in $K^{\sharp}$, hence compact in $X$, and
\[
x\in B((t,\infty];m,E)\subseteq C\subseteq B(J;m,E)\subseteq U.
\]
In both cases $C$ is upper for the specialization order, so it is
saturated. Indeed, in any space the saturation of a set $A$ is
$\up A$: if $y\notin\up A$, the union of open sets containing
individual points of $A$ and excluding $y$ is an open superset
of $A$ excluding $y$.
\end{proof}

Recall that a nonempty subset $D$ of a poset is \emph{directed}
if every two of its elements have an upper bound in $D$.
In a complete lattice, $a\ll b$ means that whenever
$b\leq\bigvee D$ for a directed set $D$, some $d\in D$ satisfies
$a\leq d$. The lattice is \emph{continuous} if
$b=\bigvee\{a:a\ll b\}$ for every $b$. A \emph{frame} is a
complete lattice in which finite meets distribute over arbitrary joins.

\begin{proposition}\label{prop:frame}
The complete lattice $\OO(X)$ is a continuous frame.
\end{proposition}

\begin{proof}
For every space, unions and finite intersections of opens obey
the frame identity. Let $x\in U\in\OO(X)$, and choose $C$ as in
Proposition~\ref{prop:compact}. Put $V=\intr_X C$.
If $\mathcal D\subseteq\OO(X)$ is directed and
$U\subseteq\bigcup\mathcal D$, compactness of $C$ gives finitely
many members of $\mathcal D$ covering $C$. A common upper bound
$D\in\mathcal D$ contains $C$, and hence $V$.
Thus $V\ll U$. Varying $x$ shows that $U$ is the union of opens
way below it. The assertion for $U=\varnothing$ follows from
$\varnothing\ll\varnothing$.
\end{proof}

\begin{proposition}\label{prop:continuous-meet}
The map $\wedge:X\times X\longrightarrow X$ in \eqref{eq:meet}
is continuous.
\end{proposition}

\begin{proof}
Let $x=(s,(A_i))$, $y=(t,(B_i))$, and suppose
$x\wedge y\in W=B(J;m,E)$. Write the first coordinate of
$x\wedge y$ as $u$. For $i\leq m$, we have $E_i\subseteq A_i\cap B_i$.
If $J=\RR$, then $x,y\in W$ and every meet of two points of $W$
belongs to $W$. Hence $W\times W\subseteq\wedge^{-1}(W)$.

Suppose $J=(r,\infty]$. Since $u>r$, choose finite real numbers
$\tau,\varepsilon$ such that
\begin{equation}\label{eq:margin}
r<\tau-\varepsilon<\tau<u,\qquad\varepsilon>0.
\end{equation}
Then $\tau>\varepsilon$, $s,t>\tau$, and
$\sigma_i(A_i\cap B_i)>\tau$ for every $i$.
Choose $i_0$ from Lemma~\ref{lem:slow-intersection}, and put
$q=\max\{m,i_0\}$. The sets
\begin{align*}
P&=X\cap\left((\tau,\infty]\times\prod_{i=0}^{q}\up A_i
                                      \times\prod_{i>q}S_i\right),\\
Q&=X\cap\left((\tau,\infty]\times\prod_{i=0}^{q}\up B_i
                                      \times\prod_{i>q}S_i\right)
\end{align*}
are open neighborhoods of $x$ and $y$.
Take $x'=(s',(A_i'))\in P$ and $y'=(t',(B_i'))\in Q$.
For $i\leq q$, monotonicity gives
\[
\sigma_i(A_i'\cap B_i')\geq\sigma_i(A_i\cap B_i)>\tau.
\]
For $i>q$, the defining constraints of $X$ imply
$\sigma_i(A_i')\geq s'>\tau$ and
$\sigma_i(B_i')\geq t'>\tau$. Lemma~\ref{lem:slow-intersection}
therefore gives $\sigma_i(A_i'\cap B_i')>\tau-\varepsilon$.
Consequently
\[
\inf_i\sigma_i(A_i'\cap B_i')\geq\tau-\varepsilon,
\]
so the first coordinate of $x'\wedge y'$ is at least
$\tau-\varepsilon>r$. Also $E_i\subseteq A_i'\cap B_i'$ for
$i\leq m$. Hence $x'\wedge y'\in W$, proving
$P\times Q\subseteq\wedge^{-1}(W)$.
\end{proof}

\begin{corollary}\label{cor:finite-dist}
The continuous frame $L=\OO(X)$ is $\kappa$-distributive for every
integer $\kappa\geq2$.
\end{corollary}

\begin{proof}
Apply Lemma~\ref{lem:finite} to Propositions~\ref{prop:order},
\ref{prop:meet}, and~\ref{prop:continuous-meet}, and use
Proposition~\ref{prop:frame}.
\end{proof}

\section{A countable obstruction to \texorpdfstring{$\omega$}{omega}-distributivity}
\label{sec:obstruction}

Let
\begin{equation}\label{eq:Ustar}
U_* =\{(s,(A_i))\in X:s>1\}.
\end{equation}
This is a nonempty open set, since $\tops\in U_*$.
For $n\geq2$ and $k\in N_n$, define
\begin{equation}\label{eq:z}
z_{n,k}=(n,(A_i^{n,k})),\qquad
A_i^{n,k}=\begin{cases}
N_n\setminus\{k\},&i=n,\\
N_i,&i\ne n.
\end{cases}
\end{equation}
Formula~\eqref{eq:sigma-special} gives $z_{n,k}\in X$.
The finite first coordinate $n$ is essential for this membership.

\begin{lemma}\label{lem:escape}
For every neighborhood $V$ of $\tops$, there is $n_0\geq2$ such that
$z_{n,k}\in V$ for all $n\geq n_0$ and all $k\in N_n$.
For each $n\geq2$, the finite family $\{z_{n,k}:k\in N_n\}$ has
no common lower bound in $U_*$.
\end{lemma}

\begin{proof}
Choose a basic open $B(J;m,E)$ containing $\tops$ and contained
in $V$. If $J=\RR$, every $n>m$ with $n\geq2$ suffices. If
$J=(r,\infty]$, every $n>\max\{m,r\}$ with $n\geq2$ suffices.
Indeed, all coordinates constrained by this basic open equal
$N_i$ at $z_{n,k}$, and the first coordinate is $n$.

If $y=(s,(D_i))$ is below all $z_{n,k}$, then
\[
D_n\subseteq\bigcap_{k\in N_n}(N_n\setminus\{k\})=\varnothing.
\]
The defining inequality for $X$ gives
$s\leq\sigma_n(\varnothing)=0$. Thus $y\notin U_*$.
\end{proof}

For each $n\geq2$ and $k\in N_n$, the open coordinate cylinder
\[
H_{n,k}=\{(s,(A_i))\in X:k\in A_n\}
\]
also has the description
\begin{equation}\label{eq:complement}
H_{n,k}=X\setminus\dn z_{n,k}.
\end{equation}
Indeed, if $k\notin A_n$, then $A_n\subseteq N_n\setminus\{k\}$
and $s\leq\sigma_n(A_n)\leq n$, so $(s,(A_i))\leq z_{n,k}$.
The converse follows from the $n$th coordinate comparison.
Define the following entries of $L=\OO(X)$:
\begin{equation}\label{eq:entries}
a_{n,k}=U_*\cap H_{n,k}\qquad(n\geq2,\ k\in N_n).
\end{equation}

\begin{theorem}\label{thm:countable}
The matrix \eqref{eq:entries} satisfies
\begin{align}
\bigwedge_{n\geq2}\bigvee_{k\in N_n}a_{n,k}&=U_*\ne\varnothing,
                                                    \label{eq:left}\\
\bigvee_{f\in\prod_{n\geq2}N_n}\bigwedge_{n\geq2}a_{n,f(n)}
 &=\varnothing.                                    \label{eq:right}
\end{align}
In particular, $L$ is not $\omega$-distributive.
\end{theorem}

\begin{proof}
If $(s,(A_i))\in U_*$, then $A_n\ne\varnothing$ for every $n$,
since otherwise $s\leq\sigma_n(\varnothing)=0$.
For fixed $n\geq2$, choose $k\in A_n$. This proves
$\bigcup_{k\in N_n}a_{n,k}=U_*$, and hence \eqref{eq:left}.

Fix $f\in\prod_{n\geq2}N_n$, and put
\[
b_f=\bigwedge_{n\geq2}a_{n,f(n)}
    =\intr_X\left(\bigcap_{n\geq2}a_{n,f(n)}\right).
\]
If $b_f$ were nonempty, Proposition~\ref{prop:order} would give
$\tops\in b_f$. Lemma~\ref{lem:escape}, applied to the open
neighborhood $b_f$, would then give $z_{n,f(n)}\in b_f$ for all
sufficiently large $n$. But $f(n)\notin A_n^{n,f(n)}$, so
$z_{n,f(n)}\notin a_{n,f(n)}$, contradicting
$b_f\subseteq a_{n,f(n)}$. Thus $b_f=\varnothing$ for every $f$,
which proves \eqref{eq:right}. Each $N_n$ is finite and nonempty,
so Lemma~\ref{lem:distribution} proves the last assertion.
\end{proof}

\begin{corollary}\label{cor:approxzero}
In Ern\'e's approximation notation,
$\DD_\omega U_*=\{\varnothing\}$ in $\OO(X)$.
Consequently, $L$ is not a wide coframe and is not completely
distributive.
\end{corollary}

\begin{proof}
Since $U_*\ne\varnothing$, the empty open set belongs to
$\DD_\omega U_*$. Suppose a nonempty open $V$ belongs to this set.
It contains $\tops$, so Lemma~\ref{lem:escape} gives $n\geq2$
with $z_{n,k}\in V$ for every $k\in N_n$.
The row $(a_{n,k})_{k\in N_n}$ covers $U_*$, but no entry contains
$V$, since $z_{n,k}\in V\setminus a_{n,k}$.
This contradicts \eqref{eq:approx}.
The final assertions also follow from Theorem~\ref{thm:countable}
and Lemma~\ref{lem:distribution}.
\end{proof}

\section{Web spaces and the finite-cardinal distinction}
\label{sec:web}

We finish by giving the topological conclusion directly in
Ern\'e's terminology. For a subset $A$ of a specialization poset,
write $A^{\downarrow}=\{y:y\leq a\text{ for all }a\in A\}$.

\begin{definition}[{\cite[Section~6]{Erne2009}}]\label{def:web}
A space $Y$ is a \emph{$\kappa$-web space} if, for every $x\in Y$
and every neighborhood $U$ of $x$, there is a neighborhood $V$ of
$x$ such that
\[
x\in A\subseteq V,\quad |A|<\kappa
\quad\Longrightarrow\quad A^{\downarrow}\cap U\ne\varnothing.
\]
An $\omega$-web space is called a \emph{wide web space}.
\end{definition}

\begin{proposition}\label{prop:web}
The space $X$ is a $\kappa$-web space for every integer
$\kappa\geq2$, but is not a wide web space.
\end{proposition}

\begin{proof}
Fix finite $\kappa\geq2$, $x\in X$, and a neighborhood $U$ of $x$.
Choose an open neighborhood $U'\subseteq U$.
The construction in \eqref{eq:finite-local}, with $n=\kappa-1$,
provides an open neighborhood $V$ of $x$ such that
$\mu_n(V^n)\subseteq U'$.
Every nonempty $A\subseteq V$ with $|A|\leq n$ can be listed as
an $n$-tuple by repetition. Its meet therefore belongs to $U'$,
and is a lower bound of $A$. This proves the finite-web assertion.

For failure of the wide web property, take $x=\tops$ and $U=U_*$.
Given any neighborhood $V$ of $\tops$, choose $n$ as in
Lemma~\ref{lem:escape}, and put
\[
A=\{\tops\}\cup\{z_{n,k}:k\in N_n\}.
\]
Then $A$ is finite, $\tops\in A\subseteq V$, and
$A^{\downarrow}\cap U_*=\varnothing$ by that lemma.
\end{proof}

\begin{proof}[Proof of Theorem~\ref{thm:main}]
The topological assertions are Propositions~\ref{prop:order},
\ref{prop:compact}, and~\ref{prop:continuous-meet}.
The lattice assertions are Proposition~\ref{prop:frame},
Corollary~\ref{cor:finite-dist}, and Theorem~\ref{thm:countable}.
The web-space assertions are Proposition~\ref{prop:web}.
\end{proof}

The finite-cardinal distinction also affects a statement printed
later in \cite{Erne2009}. We make the consequence explicit to avoid
relying on that statement outside its valid scope.

\begin{proposition}\label{prop:scope}
For every finite $\kappa\geq3$, the implication
\textup{(b)}$\Rightarrow$\textup{(a)} in
\cite[Proposition~6, p.~2066]{Erne2009} fails as printed.
Namely, a $\kappa$-web space whose specialization order is a
$\kappa$-meet-semilattice need not have a neighborhood base of
$\kappa$-meet-subsemilattices.
\end{proposition}

\begin{proof}
The space $X$ has a greatest element and binary meets, so it has
all finite meets. Proposition~\ref{prop:web} gives the
$\kappa$-web property. If $\tops$ had a neighborhood $N$ contained
in $U_*$ and closed under meets of fewer than $\kappa$ elements,
then $N$ would in particular be closed under binary meets.
It would therefore contain the meet of each nonempty finite
subset of $N$. Choose $n$ from Lemma~\ref{lem:escape} with
$z_{n,k}\in N$ for every $k\in N_n$. Their meet would belong to
$N\subseteq U_*$, contrary to the same lemma.
\end{proof}

\begin{remark}\label{rem:regular}
The proof of the cited implication passes from $\kappa$ to the
least regular cardinal $\bar\kappa\geq\kappa$ and, in the
nonregular case, uses the assertion that $\bar\kappa$ is the
cardinal successor of $\kappa$. This assertion concerns infinite
cardinals and does not cover finite $\kappa\geq3$.
With the union characterization of regularity specified in
\cite[Theorem~6]{Erne2009}, the least regular cardinal above such
a finite $\kappa$ is $\omega$. The $\omega$-small sets have
unbounded finite cardinalities, and the finite-web hypothesis
does not provide one neighborhood controlling all their meets.
Proposition~\ref{prop:scope} concerns only the finite-cardinal
scope of that printed implication. The proofs above use neither
that implication nor the general web-space characterization.
\end{remark}

\begin{thebibliography}{9}

\bibitem{Erne2009}
M.~Ern\'e,
\emph{Infinite distributive laws versus local connectedness and
compactness properties},
Topology Appl. \textbf{156} (2009), no.~12, 2054--2069.
\href{https://doi.org/10.1016/j.topol.2009.03.029}
{doi:10.1016/j.topol.2009.03.029}.

\bibitem{GierzEtAl2003}
G.~Gierz, K.~H.~Hofmann, K.~Keimel, J.~D.~Lawson,
M.~Mislove, and D.~S.~Scott,
\emph{Continuous Lattices and Domains},
Encyclopedia of Mathematics and its Applications, vol.~93,
Cambridge University Press, Cambridge, 2003.

\bibitem{GL2022}
J.~Goubault-Larrecq,
\emph{Topological semilattices with small semilattices},
Non-Hausdorff Topology and Domain Theory, electronic supplement,
2022.
\url{https://projects.lsv.ens-paris-saclay.fr/topology/?page_id=4939}
(accessed September~5, 2026).

\bibitem{Lawson1970}
J.~D.~Lawson,
\emph{Lattices with no interval homomorphisms},
Pacific J. Math. \textbf{32} (1970), no.~2, 459--465.
\href{https://doi.org/10.2140/pjm.1970.32.459}
{doi:10.2140/pjm.1970.32.459}.

\end{thebibliography}
\end{document}
